\documentclass{article}
\usepackage{PRIMEarxiv}
\usepackage{amsmath, amssymb, amsthm, bm, dcolumn}
\usepackage[numbers,sort&compress]{natbib}
\usepackage{graphicx}
\usepackage{multicol}
\usepackage{hyperref}
\usepackage{listings}
\usepackage{authblk}
\usepackage{wrapfig}
\usepackage[pscoord]{eso-pic}
\usepackage[fulladjust]{marginnote}
\usepackage[protrusion=true, expansion=true, tracking=false]{microtype}
\usepackage[right]{lineno}
\usepackage{setspace}
\doublespacing
\usepackage{changepage}
\usepackage[aboveskip=1pt,labelfont=bf,labelsep=period,singlelinecheck=off]{caption}
\usepackage{lastpage,fancyhdr}

\pagestyle{fancy}
\fancyhf{}
\fancyhead[L]{Preprint - PrimeAI Enhanced Template}
\fancyfoot[C]{\scriptsize Multiplicity Theory \copyright 2024 Ryan Van Gelder - Citizen Gardens}
\fancyfoot[R]{Page \thepage\ of \pageref{LastPage}}
\renewcommand{\footrule}{\hrule height 2pt \vspace{2mm}}

\begin{document}

\documentclass{article}
\usepackage{amsmath, amssymb, amsthm, bm, hyperref}
\usepackage{graphicx}
\usepackage{authblk}
\usepackage{geometry}
\geometry{a4paper, margin=1in}

\title{Multiplicity Theory and Black Hole Singularities: A Computational Tensor Network Approach}
\author{Citizen Gardens Research Team}
\affil{The Foundation of Multiplicity, info@citizengardens.org}
\date{2025}

\maketitle

\begin{abstract}
Multiplicity Theory introduces recursive multiplicative structures, prime-indexed eigenvalues, and tensor-based entanglement to reformulate Einstein’s Quantum Field Equations (QFEs). This paper explores how Multiplicity Theory resolves black hole singularities, preserves quantum information, and extends gravitational equations into the cognitive domain. We propose a self-referential tensor network model, demonstrating that thought itself emerges as an eigenvector within a multiplicative quantum Hilbert space.
\end{abstract}



\nocite{*}
\bibliographystyle{unsrt}
\bibliography{references}

\end{document}